\documentclass[a4paper]{article}     
\usepackage{amsfonts}
\usepackage{amsmath}
\usepackage{amssymb}
\usepackage{graphicx}
\usepackage{cancel}

\newcommand{\asa}{\Leftrightarrow} %als slechts als
\newcommand{\adR}{\Rightarrow} %als dan Rechts
\newcommand{\adL}{\Leftarrow}
\newcommand{\Lh}{\left(} %Links haakje
\newcommand{\Rh}{\right)}
\newcommand{\Lv}{\left[} %Links vierkanthaakje
\newcommand{\Rv}{\right]}
\newcommand{\La}{\left\{} %Links acolade
\newcommand{\Ra}{\right\}}
\newcommand{\Ras}{\right.} %Rechts acolade stelsel (omdat om stelsels te maken heb je een punt nodig
\newcommand{\pnR}{\rightarrow} %pijl naar Rechts
\newcommand{\limiet}{\lim\limits_}
\newcommand{\schrap}{\cancel}
\newcommand{\schrapb}{\bcancel} %backslash
\newcommand{\schrapk}{\xcancel} %kruis
\newcommand{\vervang}{\cancelto}
\newcommand{\intgrl}{\displaystyle\int} %integraal teken (bij het berekenen van primitieven)

\begin{document}

\noindent \large Auteur: Yoren Collier \\
\noindent \large Studierichting: Informatica\\
\noindent \large Oefeningenreeks 5A nr. 4.17\\

\medskip

\normalsize

$ BI \intgrl_0^1 \dfrac{dx}{5x^2 +5} \adR OBI = \intgrl \dfrac{1}{5x^2 +5} dx$\\

$= \intgrl \dfrac{1}{5 \Lh x^2 +1 \Rh} dx = \dfrac{1}{5} \intgrl \dfrac{1}{x^2 +1} dx = \dfrac{1}{5} \operatorname{Bgtan} \Lh x \Rh +c$\\

$BI = \Lv \dfrac{1}{5} \operatorname{Bgtan} \Lh x \Rh +c \Rv_0^1 = \Lh \dfrac{1}{5} \operatorname{Bgtan} \Lh 1 \Rh +c \Rh - \Lh \dfrac{1}{5} \operatorname{Bgtan} \Lh x \Rh +c \Rh = \dfrac{1}{5} \cdot \dfrac{\pi}{4} \schrap{+c} - \schrapb{\dfrac{1}{5} \cdot 0} \schrap{-c} =\dfrac{\pi}{20}$\\

\begin{thebibliography}{999}
%\bibitem{a} Referentie
\end{thebibliography}
\end{document} 